\section{Introduction}
\label{sec:introduction}

Large language models are rapidly replacing human annotators in computational social science, with researchers using LLM-generated labels for tasks ranging from sentiment classification to political stance detection~\cite{egami_hinck_stewart_wei2023, baumann2025}. This shift promises dramatic cost reduction, but introduces a critical validity gap: annotation errors from imperfect classifiers propagate into downstream causal estimates in ways that depend on both the error structure and the causal design. \citet{baumann2025} demonstrate that across 37 annotation tasks, 31\% of hypotheses yield incorrect conclusions with state-of-the-art LLMs, rising to 50\% for smaller models. Even classifiers achieving $F_1 > 0.93$ produce over 50\% downstream error risk.

The consequences can be catastrophic. Under $K{=}3$ non-differential exposure misclassification, 83.4\% of confusion matrices cause \emph{sign reversal} of the estimated treatment effect, where the estimated direction of the causal relationship is opposite to the truth. This failure breaks the classical binary guarantee that non-differential misclassification biases estimates ``toward the null''~\cite{brenner1993}, a guarantee that many applied researchers implicitly rely upon. The breakdown is not an edge case: it persists even when classifier accuracy exceeds 70\% (13.7\% violation rate at diagonal $\geq 0.7$, comprising 2.1\% sign reversal and 11.7\% amplification).

Why does the multi-class setting behave so differently? For binary classification, non-differential error can only attenuate the effect because the $2{\times}2$ stochastic confusion matrix preserves coefficient signs. With $K \geq 3$ categories, the confusion matrix can assign negative effective weights to categories with opposite-sign coefficients, enabling sign reversal~\cite{veirod_laake2001}. This fundamental asymmetry between binary and multi-class annotation error has received limited systematic treatment. \citet{brenner1993} identified the possibility of polytomous reversal but did not characterize it across causal designs. \citet{lee_wood_doughty2024} extended classifier error analysis to categorical variables and LLM classifiers, but covered only one DAG topology (unobserved confounding). No existing framework addresses how the same confusion matrix produces qualitatively different bias across different causal identification strategies.

Compounding this gap, \citet{kim_garg2025} show that LLM errors are highly correlated across models (60\% agreement when both err), undermining the common assumption that ensembling multiple LLMs can mitigate annotation error. This finding strengthens the case that analyzing the \emph{structure} of annotation error, rather than simply reducing its magnitude, is essential for valid downstream inference.

We address these gaps with three contributions:

\begin{enumerate}
\item \textbf{Multi-class bias taxonomy.} We derive bias expressions for seven canonical DAG topologies (Figure~\ref{fig:dag_taxonomy}) under $K$-class non-differential misclassification (closed-form for confounding, exposure, and mediation; numerically validated for all seven). We establish that five topologies maintain a structural sign-flip guarantee (with formal proofs for three topologies and Monte Carlo validation for all seven) ($\tau_{\text{sign\_flip}} = 0$), while exposure and IV topologies can produce sign reversal. This taxonomy reveals that bias direction is determined by the \emph{interaction} between the confusion matrix and the DAG topology, not by either alone.

\item \textbf{Annotation Sensitivity Values (ASVs).} We propose a unified sensitivity analysis framework with three complementary measures: magnitude ASV (E-value analogue for annotation error), significance ASV (hypothesis test filter for attenuating topologies), and sign-flip ASV (for exposure and IV). ASVs reveal a vulnerability spectrum from IV (maximally fragile, 52.6\% sign reversal; sign-flip ASV $= 1.1$--$6.7$) to M-bias (immune, ASV $= \infty$), providing practitioners with topology-aware robustness diagnostics.

\item \textbf{Topology-aware correction guidance.} We evaluate five correction methods (naive plug-in, MLA, DSL, MC-SIMEX, PPI++) across all seven topologies and show that correction effectiveness is itself topology-dependent. A Frisch--Waugh--Lovell invariance result proves that matrix-based corrections achieve exactly 0\% bias reduction on the target coefficient for four topologies, while PPI++ achieves 93.8\% bias reduction but requires gold-standard labels. We further show that collapsing $K{=}3$ annotations to binary introduces additional bias varying 147$\times$ across topologies.
\end{enumerate}
